\abstract{





% As LLM-based assistants become embedded in everyday work and study, lifelong learning is increasingly mediated by everyday human--LLM interaction. Yet it remains unclear whether these interactions elicit the engagement through which people build understanding. We analyse 128,334 naturalistic conversations from three public conversational corpora---WildChat, LMSYS Chat and ShareChat---adapting learning-science constructs to 979,974 turns through human-verified LLM-assisted annotation. Across 490,932 user turns, cognitive engagement appeared in 31.9\% and constructive engagement---the highest-level observable form in which users elaborate, test, revise or extend ideas---in 4.9\%. Constructive engagement was organized by user framing, task ecology and sustained exchange across the three corpora. On the assistant side, support was most informative when analysed not as a binary scaffolding signal, but as a situated match between support form and user activity: feedback-like and explanatory forms aligned with constructive engagement, and adjacent-turn analyses showed local, state- and form-dependent coupling rather than a demonstrated global trajectory. These findings make informal learning in everyday AI use measurable and show how AI-mediated work can preserve opportunities for people to build, test and extend understanding while solving real problems.

As LLM-based assistants become embedded in work and study, everyday human--LLM interaction is becoming a consequential site of informal learning. Learning science links deeper overt engagement to stronger learning outcomes, making engagement a useful observable process indicator in user-generated interactions when goals are not pre-specified by curricula.
Through a large-scale analysis of 128,569 naturalistic conversations, we confirm the presence of learning-oriented engagement at different depths during users' interactions with LLMs:
% We analyse 128,569 naturalistic conversations from WildChat, LMSYS Chat and ShareChat, labelling 981,470 turns with learning-science constructs via a human-reviewed LLM-assisted workflow. 
across 491,685 user turns, cognitive engagement appeared in 31.9\%, while constructive engagement, the deepest observable form of user engagement, occurred in 4.9\%. 
Our study further reveals the factors associated with these forms of engagement. We found that scaffolded support from the assistant (present in 31.7\% of assistant turns) was a consistent factor of richer constructive participation, with associations that varied by user framing, task ecology, support form, timing and prior user state. These findings make behavioural signatures of informal learning observable at scale and broaden how future AI assistance should be evaluated: from the efficiency of answer delivery toward the preservation of cognitive opportunities for users to reason, test ideas and construct understanding in the course of everyday problem-solving.

% These findings make behavioural signatures of informal learning observable at scale and broaden how AI assistance should be evaluated: moving beyond answer-delivery efficiency to the cognitive opportunities systems preserve for users to reason, test ideas and construct understanding during everyday problem-solving.


% These findings make behavioural signatures of informal learning observable at scale and broaden the evaluative of AI assistance beyond answer-delivery efficiency: not only how efficiently systems deliver answers, but whether they preserve cognitive opportunities for users to reason, test ideas and construct understanding during everyday problem-solving.

% These findings make behavioural signatures of informal learning observable at scale and broaden the evaluation of AI assistance beyond answer-delivery efficiency, foregrounding the cognitive opportunities that systems create for users to reason, test ideas and construct understanding during everyday problem-solving.

% These findings make behavioural signatures of informal learning observable at scale and broaden the evaluative question for AI assistance: not only how efficiently systems deliver answers, but whether they preserve cognitive opportunities for users to reason, test ideas and construct understanding during everyday problem-solving.

% These findings make behavioural signatures of informal learning observable at scale and suggest that AI assistance should be evaluated not only by the speed of answer delivery, but also by whether it preserves cognitive opportunities for users to reason, test ideas and construct understanding during everyday problem-solving.





% . These rates confirm that informal learning engagement is pervasive in naturalistic human–LLM interaction, yet its deepest forms remain rare. 

% Among the factors we examined, scaffolded support from the assistant (present in 31.7\% of assistant turns) was a consistent marker of richer constructive participation, but not as a uniform exposure: its association depended on user framing, task ecology, support form, timing and prior user state. Feedback-like and explanatory forms aligned most clearly with constructive engagement, and adjacent-turn analyses showed that support--engagement coupling was state- and form-dependent. These findings make behavioural signatures of informal learning observable at scale and call for a shift in how AI assistance is evaluated—from speed of answer delivery toward the preservation of cognitive opportunities for users to reason, test ideas and construct understanding in the course of everyday problem-solving.


% We analyse 128,569 naturalistic conversations from WildChat, LMSYS Chat and ShareChat, labelling 981,470 turns with learning-science constructs via a human-reviewed LLM-assisted workflow.


% These findings make behavioural signatures of informal learning observable at scale and recast AI assistance around cognitive responsibility: whether systems merely accelerate answer delivery or also preserve opportunities for people to reason, test ideas and build understanding while solving real problems.



% scaffolded support appeared in 31.7\% of assistant turns. Scaffolded support marked richer constructive participation across settings, but not as a uniform exposure: its association depended on user framing, task ecology, support form, timing and prior user state. 


% Feedback-like and explanatory forms aligned most clearly with constructive engagement, and adjacent-turn analyses showed that support--engagement coupling was state- and form-dependent rather than a conversation-wide accumulation pattern. These findings make behavioural signatures of informal learning observable at scale and recast AI assistance around cognitive responsibility: whether systems merely accelerate answer delivery or also preserve opportunities for people to reason, test ideas and build understanding while solving real problems.


% These rates confirm that informal learning engagement is pervasive in naturalistic human–LLM interaction, yet its deepest forms remain rare. Among the factors we examined, scaffolded support from the assistant (present in 31.7\% of assistant turns) was a consistent marker of richer constructive participation, but not as a uniform exposure: its association depended on user framing, task ecology, support form, timing and prior user state. 

% Learning science links deeper overt engagement to stronger learning outcomes, making engagement a useful observable process indicator in informal, user-generated interactions whose goals are not pre-specified by curricula or study protocols.

% \haotian{I feel that the abstract is a little dense as there are too many findings. It can make readers shift their attention to something that is not what we want to highlight, such as the discussion about subject, i.e., coding vs. writing, in the sentence starting with ``However.'' What about removing the discussion and highlighting the findings in Sec. 2.4 (maybe we can call them ``the mechanism/factor of learning-oriented engagements'')?}

}
